% Native LaTeX visuals for neurodivergent-friendly chunking and source connection.
\definecolor{SoftGreen}{HTML}{E5F2E8}
\definecolor{SoftGold}{HTML}{FFF2CC}
\definecolor{SoftLavender}{HTML}{EEE8F6}
\definecolor{DeepBlue}{HTML}{244A73}
\definecolor{DeepGreen}{HTML}{315C45}
\definecolor{DeepGold}{HTML}{705B00}
\definecolor{DeepPurple}{HTML}{51416E}

\captionsetup{font=small,labelfont=bf}

\newtcolorbox{visualbox}[2][]{
  enhanced,
  breakable,
  colback=SoftGray,
  colframe=PrimaryRed,
  boxrule=0.8pt,
  arc=2mm,
  left=2mm,right=2mm,top=1.5mm,bottom=1.5mm,
  title={#2},
  fonttitle=\bfseries,
  coltitle=black,
  attach boxed title to top left={xshift=3mm,yshift=-2mm},
  boxed title style={colback=SoftBlue,colframe=PrimaryRed,boxrule=0.6pt,arc=1.5mm},
  #1
}

\newcommand{\actioncard}[3]{%
\begin{tcolorbox}[colback=#2,colframe=#3,boxrule=0.7pt,arc=2mm,
  left=1.5mm,right=1.5mm,top=1mm,bottom=1mm,height=1.05in,valign=center]
\centering\textbf{#1}
\end{tcolorbox}%
}

\newcommand{\SourceMapVisual}{%
\begin{figure}[ht]
\centering
\resizebox{0.96\textwidth}{!}{%
\begin{tikzpicture}[
  node distance=7mm and 10mm,
  source/.style={rounded corners=2mm, draw=DeepBlue, fill=SoftBlue, minimum width=2.55in, minimum height=0.68in, align=center, font=\small},
  evidence/.style={rounded corners=2mm, draw=DeepGreen, fill=SoftGreen, minimum width=2.55in, minimum height=0.68in, align=center, font=\small},
  process/.style={rounded corners=2mm, draw=PrimaryRed, fill=SoftGray, minimum width=2.75in, minimum height=0.7in, align=center, font=\small},
  outcome/.style={rounded corners=2mm, draw=DeepPurple, fill=SoftLavender, minimum width=3.15in, minimum height=0.72in, align=center, font=\small},
  arrow/.style={-{Latex[length=2.5mm]}, thick, draw=black!65}
]
\node[source] (oar) {\textbf{OAR educator resources}\\peer-friendly autism acceptance and friendship actions};
\node[evidence, right=of oar] (placement) {\textbf{Placement evidence}\\login friction, noise, pacing, help, and belonging};
\node[process, below=of $(oar)!0.5!(placement)$] (lesson) {\textbf{10-minute lesson design}\\translate ideas into concrete classroom choices};
\node[outcome, below=of lesson] (future) {\textbf{Reflection and future practice}\\build access, dignity, agency, and peer belonging};
\draw[arrow] (oar) -- (lesson);
\draw[arrow] (placement) -- (lesson);
\draw[arrow] (lesson) -- (future);
\end{tikzpicture}%
}
\caption{Visual source map showing how OAR resources and non-identifying placement evidence informed the lesson, reflection, and future classroom plan.}
\label{fig:source-map}
\end{figure}%
}

\newcommand{\LessonTimelineVisual}{%
\begin{visualbox}{10-Minute Lesson at a Glance}
\renewcommand{\arraystretch}{1.35}
\begin{tabularx}{\textwidth}{>{\bfseries\centering\arraybackslash}p{0.85in} X X}
\toprule
Time & Teaching move & Purpose \\
\midrule
0--2 min & Preview first/next/last and introduce autism as one way a brain can work. & Reduce uncertainty and avoid a single story about autism. \\
2--4 min & Name different communication and sensory needs. & Normalize speech, gestures, writing, AAC, movement, silence, space, and wait time. \\
4--8 min & Teach and practice five peer actions with a technology-class scenario. & Move from general kindness to observable supportive behavior. \\
8--10 min & Complete an exit response: ``One way I can be a supportive classmate is... because...'' & Check whether the student can connect an action to access or belonging. \\
\bottomrule
\end{tabularx}
\end{visualbox}%
}

\newcommand{\PeerActionsVisual}{%
\begin{figure}[ht]
\centering
\begin{tabularx}{\textwidth}{*{5}{>{\centering\arraybackslash}X}}
\actioncard{INVITE\\without forcing}{SoftBlue}{DeepBlue} &
\actioncard{WAIT\\give time to respond or try}{SoftGreen}{DeepGreen} &
\actioncard{ASK\\instead of guessing}{SoftGold}{DeepGold} &
\actioncard{RESPECT\\space, tools, breaks, privacy}{SoftLavender}{DeepPurple} &
\actioncard{INCLUDE\\different ways of participating}{SoftGray}{PrimaryRed}
\end{tabularx}
\caption{Five peer actions adapted from OAR friendship guidance: invite, wait, ask, respect, and include.}
\label{fig:peer-actions}
\end{figure}%
}

\newcommand{\EvidenceVisual}{%
\begin{visualbox}{From Surface Behavior to Richer Evidence}
\renewcommand{\arraystretch}{1.35}
\begin{tabularx}{\textwidth}{>{\bfseries\RaggedRight\arraybackslash}p{1.55in} X X}
\toprule
What peers may notice & A richer interpretation & A supportive response \\
\midrule
A delayed answer & The student may need processing time or another communication mode. & Wait, offer choices, and notice gestures, writing, AAC, movement, or silence. \\
Moving away or covering ears & The noise, crowding, or sensory load may be too high. & Lower the volume, create space, and ask before following or touching. \\
Difficulty joining a coding task & Login steps, reading load, peer pacing, or unclear directions may be hiding competence. & Show one step, keep the device with the learner, and let the learner try. \\
A new or quiet student stays outside the group & Belonging may not yet be socially established. & Invite without pressure and explicitly treat the person as part of the class. \\
\bottomrule
\end{tabularx}
\end{visualbox}%
}

\newcommand{\PuzzleEquitasVisual}{%
\begin{figure}[ht]
\centering
\resizebox{0.96\textwidth}{!}{%
\begin{tikzpicture}[
  piece/.style={rounded corners=2mm, minimum width=1.4in, minimum height=0.65in, align=center, font=\small\bfseries, draw=black!55},
  lens/.style={rounded corners=2mm, minimum width=2.25in, minimum height=0.92in, align=center, font=\small, draw=PrimaryRed, fill=SoftGray},
  goal/.style={rounded corners=2mm, minimum width=3.25in, minimum height=0.8in, align=center, font=\small\bfseries, draw=DeepPurple, fill=SoftLavender},
  arrow/.style={-{Latex[length=2.5mm]}, thick, draw=black!65}
]
\node[piece, fill=SoftBlue] (lived) {Lived\\experience};
\node[piece, fill=SoftGreen, right=5mm of lived] (education) {Educational\\practice};
\node[piece, fill=SoftGold, right=5mm of education] (data) {Classroom\\evidence};
\node[piece, fill=SoftLavender, right=5mm of data] (research) {Research and\\resources};
\node[lens, below=9mm of $(education)!0.5!(data)$] (equitas) {\textbf{EQUITAS equity questions}\\Whose knowledge is included?\\What context is missing? Where may bias shape the decision?};
\node[goal, below=8mm of equitas] (goal) {Access $\rightarrow$ Dignity $\rightarrow$ Agency $\rightarrow$ Belonging};
\draw[arrow] (lived) -- (equitas);
\draw[arrow] (education) -- (equitas);
\draw[arrow] (data) -- (equitas);
\draw[arrow] (research) -- (equitas);
\draw[arrow] (equitas) -- (goal);
\node[draw=PrimaryRed, dashed, rounded corners=3mm, fit=(lived)(education)(data)(research), inner sep=3mm, label={[font=\small\bfseries,text=PrimaryRed]above:Puzzle Plan: connect the pieces rather than isolate one explanation}] {};
\end{tikzpicture}%
}
\caption{Puzzle Plan and EQUITAS visual: multiple evidence sources are integrated, checked for missing context and bias, and translated into access-to-agency goals.}
\label{fig:puzzle-equitas}
\end{figure}%
}

\newcommand{\FutureRoutineVisual}{%
\begin{visualbox}{Future Classroom Routine: Action Plus Access Support}
\renewcommand{\arraystretch}{1.28}
\begin{tabularx}{\textwidth}{>{\bfseries}p{0.85in} X X}
\toprule
Action & What peers practice & What the classroom provides \\
\midrule
Invite & Ask someone to join without pressuring them. & Multiple entry points and permission to observe before joining. \\
Wait & Give classmates time to process, respond, or try. & Predictable pauses, visual directions, and reduced rushing. \\
Ask & Use respectful questions instead of assumptions. & Communication choices, help cards, and modeled consent language. \\
Respect & Honor space, sensory needs, tools, breaks, and privacy. & Quiet options, sensory flexibility, and private support pathways. \\
Include & Treat different participation as real participation. & Flexible roles, multiple response modes, and explicit belonging language. \\
\bottomrule
\end{tabularx}
\end{visualbox}%
}
